Notación: \(\band\) AND, \(\bor\) OR, \(\bxor\) XOR (\texttt{\textasciicircum{}} en Python), \(\bnot\) NOT. En Python los enteros se comportan como si tuvieran infinitos bits (complemento a dos), por lo que \(-1\) equivale a ``todos los bits en 1''. Las identidades entre AND, OR, XOR y NOT valen también para negativos.

\textbf{AND, OR y XOR.}
\begin{props}
  \item \textbf{Conmutativas:} \(a \band b = b \band a\), y lo mismo para \(\bor\) y \(\bxor\).
  \item \textbf{Asociativas:} \((a \band b) \band c = a \band (b \band c)\), y lo mismo para \(\bor\) y \(\bxor\).
  \item \textbf{Elemento neutro:} \(a \band (-1) = a\), \(\ a \bor 0 = a\), \(\ a \bxor 0 = a\).
  \item \textbf{Absorbente:} \(a \band 0 = 0\), \(\ a \bor (-1) = -1\).
  \item \textbf{Idempotencia:} \(a \band a = a\), \(\ a \bor a = a\). En cambio \(a \bxor a = 0\).
  \item \textbf{Absorción:} \(a \band (a \bor b) = a\), \(\ a \bor (a \band b) = a\).
\end{props}

\textbf{Distributivas.}
\begin{props}
  \item \(a \band (b \bor c) = (a \band b) \bor (a \band c)\)
  \item \(a \bor (b \band c) = (a \bor b) \band (a \bor c)\)
  \item \(a \band (b \bxor c) = (a \band b) \bxor (a \band c)\) \quad (AND distribuye sobre XOR)
\end{props}

\textbf{NOT.}
\begin{props}
  \item \textbf{De Morgan:} \(\bnot(a \band b) = \bnot a \bor \bnot b\), \(\ \bnot(a \bor b) = \bnot a \band \bnot b\).
  \item \(\bnot \bnot a = a\), \(\quad \bnot a = -a - 1\), \(\quad -a = \bnot a + 1 = \bnot(a - 1)\).
  \item \(a \bxor (-1) = \bnot a\) (voltea todos los bits). Con una máscara de \(k\) bits: \(a \bxor (2^k - 1) = (2^k - 1) - a\) para \(0 \le a < 2^k\).
\end{props}

\textbf{Propiedades específicas de XOR.}
\begin{props}
  \item \textbf{Autoinversa:} \((a \bxor b) \bxor b = a\). Si \(a \bxor b = c\), entonces \(a \bxor c = b\) y \(b \bxor c = a\) (permite ``despejar'' un valor).
  \item \(a \bxor b = 0 \iff a = b\).
  \item \textbf{Intercambio sin variable temporal:} \(a \gets a \bxor b;\ \ b \gets a \bxor b;\ \ a \gets a \bxor b\).
  \item XOR es la suma módulo 2 de cada bit. Por eso la paridad de los bits encendidos cumple \(\operatorname{pc}(a \bxor b) \equiv \operatorname{pc}(a) + \operatorname{pc}(b) \pmod 2\), donde \(\operatorname{pc}\) es la cantidad de bits en 1.
  \item \textbf{XOR acumulado} \(0 \bxor 1 \bxor \cdots \bxor n\): vale \(n\) si \(n \bmod 4 = 0\); \(1\) si \(n \bmod 4 = 1\); \(n+1\) si \(n \bmod 4 = 2\); \(0\) si \(n \bmod 4 = 3\). El XOR del rango \([l, r]\) es entonces (XOR hasta \(r\)) \(\bxor\) (XOR hasta \(l-1\)).
\end{props}
